\documentclass[12pt,a4paper]{article}
\usepackage[utf8]{inputenc}
\usepackage[T1]{fontenc}
\usepackage[spanish]{babel}
\usepackage{geometry}
\usepackage{graphicx}
\usepackage{hyperref}
\usepackage{xcolor}
\usepackage{enumitem}
\usepackage{booktabs}
\usepackage{array}
\usepackage{fancyhdr}
\usepackage{titlesec}
\usepackage{tcolorbox}
\usepackage{mdframed}
\usepackage{amsmath}

\geometry{margin=2.5cm}

% Colores corporativos
\definecolor{accessred}{RGB}{175,0,0}
\definecolor{lightgray}{RGB}{245,245,245}
\definecolor{darkblue}{RGB}{0,70,127}
\definecolor{noteblue}{RGB}{230,242,255}
\definecolor{warnyellow}{RGB}{255,248,220}

% Estilo de cabecera
\pagestyle{fancy}
\fancyhf{}
\fancyhead[L]{\textcolor{accessred}{\textbf{UT6 -- Bases de Datos Ofimáticas con Microsoft Access}}}
\fancyhead[R]{\textcolor{gray}{\thepage}}
\fancyfoot[C]{\textcolor{gray}{\small AOF -- Natalia Marcos Iglesias}}
\renewcommand{\headrulewidth}{1.5pt}
\renewcommand{\headrule}{\hbox to\headwidth{\color{accessred}\leaders\hrule height \headrulewidth\hfill}}

% Estilo de secciones
\titleformat{\section}{\Large\bfseries\color{accessred}}{}{0em}{\thesection.\quad}[\titlerule]
\titleformat{\subsection}{\large\bfseries\color{darkblue}}{}{0em}{\thesubsection.\quad}
\titleformat{\subsubsection}{\normalsize\bfseries\color{darkblue}}{}{0em}{\thesubsubsection.\quad}

% Cajas de nota y aviso
\tcbuselibrary{skins,breakable}

\newtcolorbox{nota}{
  colback=noteblue, colframe=darkblue,
  title={\textbf{Nota}}, fonttitle=\bfseries,
  breakable, enhanced
}

\newtcolorbox{importante}{
  colback=warnyellow, colframe=accessred,
  title={\textbf{Importante}}, fonttitle=\bfseries,
  breakable, enhanced
}

\newtcolorbox{definicion}{
  colback=lightgray, colframe=gray!60,
  title={\textbf{Definición}}, fonttitle=\bfseries,
  breakable, enhanced
}

% Estilo de código/ruta
\newcommand{\ruta}[1]{\texttt{\small\color{darkblue}#1}}
\newcommand{\tecla}[1]{\fbox{\small\texttt{#1}}}
\newcommand{\menu}[1]{\textit{\textcolor{darkblue}{#1}}}

% Título
\title{
  \vspace{-1cm}
  {\color{accessred}\rule{\linewidth}{3pt}}\\[0.3cm]
  {\Huge\bfseries\color{accessred} UT6: Utilización de Bases de\\[0.2cm]
  Datos Ofimáticas}\\[0.3cm]
  {\Large\color{darkblue} Microsoft Access}\\[0.3cm]
  {\color{accessred}\rule{\linewidth}{1pt}}\\[0.2cm]
  {\large\color{gray} Módulo: Aplicaciones Ofimáticas (AOF)}
}
\author{\textbf{Natalia Marcos Iglesias}}
\date{}

\begin{document}

\maketitle
\thispagestyle{empty}

\tableofcontents
\newpage

% =============================================
\section{Conceptos básicos de bases de datos}
% =============================================

\subsection{¿Qué es una base de datos?}

\begin{definicion}
Una \textbf{base de datos} es un conjunto de datos organizados para un uso determinado. El conjunto de programas que permite gestionar esos datos se denomina \textbf{Sistema Gestor de Bases de Datos (SGBD)}.
\end{definicion}

Las bases de datos de Access utilizan la extensión \texttt{.ACCDB}. La mayoría de los SGBD modernos emplean el \textbf{modelo relacional}, en el que los datos se organizan en \textbf{tablas}.

\subsection{Tablas de datos}

Una \textbf{tabla} almacena información sobre un tema concreto (clientes, pedidos, productos...). Está formada por:

\begin{itemize}
  \item \textbf{Columnas/Campos}: almacenan un tipo de dato concreto (nombre, dirección, código...).
  \item \textbf{Filas/Registros}: conjunto de todos los campos de un mismo elemento.
\end{itemize}

\subsubsection{Clave principal (Primary Key)}

\begin{importante}
Una \textbf{clave principal} es una columna (o combinación de columnas) que identifica de forma \emph{única} cada fila. No puede haber valores duplicados ni valores nulos en ella.
\end{importante}

\subsubsection{Clave foránea (Foreign Key)}

Una \textbf{clave foránea} es una columna que hace referencia a la clave principal de otra tabla. Garantiza la \textbf{integridad referencial}: los datos referenciados deben existir en la tabla principal.

\begin{nota}
Una tabla tiene \textbf{una única clave primaria} y puede tener \textbf{cero o más claves foráneas}.
\end{nota}

\subsubsection{Valor NULL}

El valor \texttt{NULL} indica ausencia de dato. \textbf{No es lo mismo que el valor 0 o una cadena vacía}.

\subsection{Objetos de una base de datos Access}

\begin{center}
\renewcommand{\arraystretch}{1.4}
\begin{tabular}{>{\bfseries\color{accessred}}p{3cm} p{10cm}}
\toprule
Objeto & Descripción \\
\midrule
Tabla & Almacena los datos organizados en filas y columnas. \\
Consulta & Extrae, filtra o modifica datos de una o varias tablas. \\
Formulario & Interfaz gráfica para introducir y visualizar datos. \\
Informe & Presenta datos formateados para impresión o exportación. \\
Macro & Automatiza acciones en respuesta a eventos. \\
Módulo & Procedimientos personalizados en Visual Basic (VBA). \\
\bottomrule
\end{tabular}
\end{center}

% =============================================
\section{Interfaz de Microsoft Access}
% =============================================

\subsection{Pantalla inicial}

Al abrir Access aparece la \textbf{vista Backstage} (pestaña Archivo) con:
\begin{itemize}
  \item Crear una nueva base de datos en blanco o desde plantilla.
  \item Lista de bases de datos recientes.
  \item Acceso a OneDrive para guardar en la nube.
\end{itemize}

\subsection{Elementos del entorno de trabajo}

Una vez abierta una base de datos, el entorno de trabajo incluye:

\begin{description}[leftmargin=3cm, style=nextline]
  \item[Barra de título] Muestra el nombre del archivo. Incluye botones de minimizar, maximizar y cerrar.
  \item[Barra de acceso rápido] Accesos directos personalizables: Guardar (\tecla{Ctrl+G}), Deshacer (\tecla{Ctrl+Z}), Rehacer.
  \item[Cinta de opciones] Organizada en \textbf{pestañas} $\rightarrow$ \textbf{grupos} $\rightarrow$ \textbf{botones}. Notación: \menu{INICIO > Portapapeles > Pegar}.
  \item[Panel de navegación] Muestra y gestiona todos los objetos de la base de datos.
  \item[Área de trabajo] Zona central donde se abren los objetos.
  \item[Barra de estado] Parte inferior; indica el estado actual y permite cambiar de vista.
\end{description}

\begin{nota}
Pulsando \tecla{ALT} se activa el modo de acceso por teclado: aparecen letras sobre cada botón para acceder sin ratón.
\end{nota}

\subsection{La pestaña Archivo (vista Backstage)}

La pestaña \textbf{ARCHIVO} abre la \textbf{vista Backstage}, con fondo granate, donde se gestionan los archivos:
\begin{itemize}
  \item \menu{Información}: Compactar/reparar y cifrar con contraseña.
  \item \menu{Nuevo}: Crear nueva base de datos.
  \item \menu{Abrir}: Abrir base de datos existente.
  \item \menu{Guardar}: Guardar cambios.
  \item \menu{Imprimir}: Opciones de impresión.
  \item \menu{Cerrar}: Cerrar la base de datos sin salir de Access.
\end{itemize}

% =============================================
\section{Gestión de bases de datos}
% =============================================

\subsection{Crear una base de datos}

\begin{enumerate}
  \item \menu{Archivo > Nuevo > Base de datos del escritorio en blanco}.
  \item Escribir el nombre del archivo (extensión \texttt{.accdb}).
  \item Seleccionar la ubicación haciendo clic en el icono de carpeta.
  \item Pulsar \textbf{Crear}.
\end{enumerate}

\begin{nota}
Al abrir, Access puede mostrar una advertencia sobre macros VBA. Si el archivo es de confianza, pulsar \textbf{Habilitar contenido}.
\end{nota}

\subsection{Abrir una base de datos existente}

\begin{itemize}
  \item Desde la lista de \textbf{Recientes} al abrir Access.
  \item \menu{Archivo > Abrir} (o \tecla{Ctrl+A}).
  \item Doble clic sobre el archivo \texttt{.accdb} en el explorador.
\end{itemize}

\subsection{Cerrar una base de datos}

\begin{itemize}
  \item \menu{Archivo > Cerrar} (mantiene Access abierto).
  \item Cerrar Access directamente cierra también la base de datos.
\end{itemize}

% =============================================
\section{Panel de navegación}
% =============================================

El \textbf{Panel de navegación} (lateral izquierdo) organiza todos los objetos de la base de datos. Se puede:

\begin{itemize}
  \item \textbf{Ocultar/mostrar}: botón de flecha en su parte superior.
  \item \textbf{Cambiar agrupación}: clic en la cabecera del panel para elegir entre \emph{Tipo de objeto}, \emph{Fecha de creación/modificación}, \emph{Tablas y vistas relacionadas} o \emph{Personalizado}.
  \item \textbf{Filtrar}: seleccionar qué tipo de objetos mostrar.
  \item \textbf{Ocultar objetos individualmente}: clic derecho $\rightarrow$ \menu{Ocultar en este grupo}.
\end{itemize}

% =============================================
\section{Creación y gestión de tablas}
% =============================================

\subsection{Crear una tabla en Vista Diseño}

Para crear una tabla:
\begin{enumerate}
  \item \menu{Crear > Tablas > Vista Diseño}.
  \item Definir cada campo indicando:
    \begin{itemize}
      \item \textbf{Nombre del campo}: identificador único en la tabla.
      \item \textbf{Tipo de datos}: determina qué valores acepta.
      \item \textbf{Descripción} (opcional): texto explicativo que aparece en la barra de estado.
    \end{itemize}
  \item Asignar la \textbf{clave principal}.
  \item Guardar la tabla con un nombre.
\end{enumerate}

\subsection{Tipos de datos}

\begin{center}
\renewcommand{\arraystretch}{1.3}
\begin{tabular}{>{\bfseries\color{darkblue}}p{3.5cm} p{9cm}}
\toprule
Tipo & Descripción \\
\midrule
Texto corto & Hasta 255 caracteres. Para nombres, teléfonos, códigos. \\
Texto largo & Más de 255 caracteres. Comentarios, descripciones largas. \\
Número & Valores numéricos para cálculos (Byte, Entero, Entero largo, Simple, Doble, Decimal). \\
Fecha/Hora & Fechas y horas (años 100--9999). \\
Moneda & Valores monetarios con alta precisión (hasta 15 dígitos enteros y 4 decimales). \\
Autonumeración & Entero largo generado automáticamente por Access. No editable. \\
Sí/No & Valores booleanos: Sí/No, Verdadero/Falso, Activado/Desactivado. \\
Objeto OLE & Imágenes, hojas de cálculo, documentos incrustados. Obsoleto desde Access 2007. \\
Hipervínculo & URL o ruta de archivo. \\
Datos adjuntos & Archivos adjuntos (imágenes, PDFs...). Más eficiente que OLE. \\
Calculado & Campo cuyo valor resulta de una expresión sobre otros campos. \\
Asistente búsquedas & Lista de valores o referencia a otra tabla (cuadro combinado). \\
\bottomrule
\end{tabular}
\end{center}

\subsection{Clave principal}

Para asignar la clave principal:
\begin{enumerate}
  \item Hacer clic sobre el campo deseado en Vista Diseño.
  \item \menu{Diseño > Herramientas > Clave principal}.
\end{enumerate}

Aparecerá una llave (\textbf{$\mathbf{\keyindicator}$}) a la izquierda del campo.

\begin{importante}
Para una \textbf{clave compuesta} (varios campos): seleccionar todos los campos con \tecla{Ctrl}+clic y luego pulsar \menu{Clave principal}.
\end{importante}

\subsection{Propiedades de los campos}

En Vista Diseño, la zona inferior muestra las propiedades del campo seleccionado:

\subsubsection{Tamaño del campo}
\begin{itemize}
  \item \textbf{Texto corto}: máximo 255 caracteres (por defecto 255).
  \item \textbf{Número}: Byte (0--255), Entero (±32.767), Entero largo (±2.147.483.647), Simple, Doble, Decimal.
\end{itemize}

\subsubsection{Formato}
Personaliza cómo se muestran los datos. Formatos numéricos: Número general, Moneda, Euro, Fijo, Estándar, Porcentaje, Científico. Formatos fecha: Fecha general, Fecha larga, Fecha mediana, Fecha corta, Hora larga, Hora mediana, Hora corta.

\subsubsection{Máscara de entrada}
Plantilla que guía la introducción de datos. Ejemplo: \texttt{(\_\_\_) \_\_\_-\_\_\_\_} para teléfonos. Solo funciona con tipos Texto y Fecha.

\subsubsection{Título}
Texto que aparece como encabezado de columna en lugar del nombre del campo.

\subsubsection{Valor predeterminado}
Valor que se introduce automáticamente si el usuario no escribe nada.

\subsubsection{Regla de validación}
Condición que debe cumplir el valor introducido. Ejemplo: \texttt{>=100 Y <=2000}.

\subsubsection{Texto de validación}
Mensaje de error que aparece cuando no se cumple la regla de validación.

\subsubsection{Requerido}
Si es \texttt{Sí}, el campo no puede quedar vacío.

\subsubsection{Indexado}
Mejora la velocidad de búsquedas y ordenaciones. Opciones: \emph{No}, \emph{Sí (con duplicados)}, \emph{Sí (sin duplicados)}.

\subsection{Trabajar en Vista Hoja de datos}

La Vista Hoja de datos permite introducir, modificar y eliminar registros:
\begin{itemize}
  \item \textbf{Introducir}: escribir en la primera fila vacía y pulsar \tecla{Intro} para avanzar de campo.
  \item \textbf{Eliminar registro}: seleccionar fila $\rightarrow$ \menu{Inicio > Registros > Eliminar} o tecla \tecla{Supr}.
  \item \textbf{Desplazarse}: barra de navegación inferior o teclado (\tecla{$\uparrow$}/\tecla{$\downarrow$} entre registros, \tecla{$\leftarrow$}/\tecla{$\rightarrow$} entre campos).
  \item \textbf{Buscar}: \menu{Inicio > Buscar} o \tecla{Ctrl+B}.
\end{itemize}

% =============================================
\section{Relaciones entre tablas}
% =============================================

\subsection{Tipos de relaciones}

\begin{center}
\renewcommand{\arraystretch}{1.4}
\begin{tabular}{>{\bfseries\color{accessred}}p{3.5cm} p{9cm}}
\toprule
Tipo & Descripción y ejemplo \\
\midrule
Uno a uno (1:1) & Cada registro de A se relaciona con uno de B y viceversa. \newline\textit{Ej.: Población -- Alcalde.} \\
Uno a varios (1:N) & Un registro de A puede relacionarse con varios de B, pero cada registro de B solo con uno de A. \newline\textit{Ej.: Población -- Habitantes.} \\
Varios a varios (N:M) & Un registro de A puede relacionarse con varios de B y viceversa. Requiere \textbf{tabla intermedia}. \newline\textit{Ej.: Clientes -- Artículos (tabla Líneas de pedido).} \\
\bottomrule
\end{tabular}
\end{center}

\subsection{Crear una relación}

\begin{enumerate}
  \item \menu{Herramientas de base de datos > Relaciones}.
  \item Agregar las tablas desde \menu{Mostrar tabla}.
  \item Arrastrar el campo clave de la tabla principal hasta el campo correspondiente de la tabla secundaria.
  \item En el cuadro \textbf{Modificar relaciones}: verificar campos, activar \textbf{Exigir integridad referencial}.
  \item Pulsar \textbf{Crear}.
\end{enumerate}

\subsection{Integridad referencial}

Al activar la integridad referencial:
\begin{itemize}
  \item No se puede insertar en la tabla secundaria un valor que no exista en la principal.
  \item No se puede eliminar de la principal un registro que tenga relacionados en la secundaria (salvo con las opciones en cascada).
\end{itemize}

\textbf{Opciones en cascada:}
\begin{description}
  \item[Actualizar en cascada] Al cambiar la clave principal, se actualizan automáticamente todas las claves foráneas relacionadas.
  \item[Eliminar en cascada] Al borrar un registro principal, se borran automáticamente todos los registros secundarios relacionados.
\end{description}

\begin{importante}
El campo de unión debe ser del \textbf{mismo tipo de dato} en ambas tablas, aunque puede tener distinto nombre.
\end{importante}

\subsection{Gestionar relaciones}

\begin{itemize}
  \item \textbf{Modificar}: doble clic sobre la línea de relación.
  \item \textbf{Eliminar}: seleccionar la línea $\rightarrow$ \tecla{Supr}.
  \item \textbf{Agregar tabla}: \menu{Diseño > Mostrar tabla}.
  \item \textbf{Ocultar tabla}: clic derecho $\rightarrow$ \menu{Ocultar tabla}.
  \item \textbf{Limpiar ventana}: \menu{Diseño > Borrar diseño} (no elimina las relaciones, solo las oculta).
  \item \textbf{Ver relaciones directas}: clic derecho sobre tabla $\rightarrow$ \menu{Mostrar directas}.
\end{itemize}

% =============================================
\section{Consultas}
% =============================================

\subsection{Tipos de consultas}

\begin{description}
  \item[Consultas de selección] Extraen y muestran datos que cumplen ciertos criterios. Generan una tabla lógica (en memoria).
  \item[Consultas de acción] Realizan cambios en los datos: creación de tabla, actualización, datos anexados, eliminación.
  \item[Consultas SQL] Se definen directamente en lenguaje SQL (consultas de unión, etc.). No se estudian en este módulo.
\end{description}

\subsection{Crear una consulta en Vista Diseño}

\begin{enumerate}
  \item \menu{Crear > Consultas > Diseño de Consulta}.
  \item Seleccionar las tablas de origen y pulsar \textbf{Agregar} $\rightarrow$ \textbf{Cerrar}.
  \item En la \textbf{cuadrícula QBE} (zona inferior), agregar campos y definir la consulta.
  \item Guardar con \tecla{Ctrl+G} o \menu{Archivo > Guardar}.
  \item Ejecutar con el botón \textbf{Ejecutar} (\menu{Diseño}) o desde el Panel de Navegación (doble clic).
\end{enumerate}

\subsection{La cuadrícula QBE}

\begin{center}
\renewcommand{\arraystretch}{1.4}
\begin{tabular}{>{\bfseries\color{darkblue}}p{2.5cm} p{10cm}}
\toprule
Fila & Función \\
\midrule
Campo & Nombre del campo o expresión calculada. \\
Tabla & Tabla de procedencia del campo. \\
Orden & Tipo de ordenación: Ascendente / Descendente. \\
Mostrar & Casilla: si está activada, el campo aparece en el resultado. \\
Criterios & Condición que deben cumplir los registros. \\
O & Condiciones alternativas (unidas por OR). \\
\bottomrule
\end{tabular}
\end{center}

\subsection{Añadir campos}

\begin{itemize}
  \item Doble clic sobre el campo en la zona de tablas.
  \item Arrastrar el campo desde la tabla hasta la cuadrícula.
  \item Hacer clic en la fila Campo de una columna vacía y seleccionar de la lista.
  \item Usar \texttt{*} para incluir todos los campos de una tabla.
\end{itemize}

\subsection{Campos calculados}

Se definen con la sintaxis: \texttt{NombreColumna: expresión}

Ejemplo: \texttt{Precio con IVA: [Precio]*1,21}

\subsubsection{Operadores disponibles}

\begin{center}
\begin{tabular}{c l}
\toprule
Operador & Uso \\
\midrule
\texttt{+} & Suma numérica \\
\texttt{-} & Resta \\
\texttt{*} & Multiplicación \\
\texttt{/} & División (resultado decimal) \\
\texttt{\textbackslash} & División entera \\
\texttt{\^{}} & Potencia \\
\texttt{Mod} & Resto de la división \\
\texttt{\&} & Concatenación de texto \\
\bottomrule
\end{tabular}
\end{center}

\subsubsection{Funciones predefinidas útiles}

\begin{itemize}
  \item \texttt{Date()} / \texttt{Fecha()} --- Fecha actual.
  \item \texttt{Now()} / \texttt{Hoy()} --- Fecha y hora actual.
  \item \texttt{Year(fecha)} / \texttt{Año()} --- Año de una fecha.
  \item \texttt{Month(fecha)} / \texttt{Mes()} --- Mes de una fecha.
\end{itemize}

\subsection{Criterios de búsqueda}

\subsubsection{Operadores de comparación}

\begin{center}
\begin{tabular}{c l}
\toprule
Operador & Significado \\
\midrule
\texttt{=} & Igual \\
\texttt{<>} & Distinto \\
\texttt{<} & Menor que \\
\texttt{<=} & Menor o igual \\
\texttt{>} & Mayor que \\
\texttt{>=} & Mayor o igual \\
\bottomrule
\end{tabular}
\end{center}

\subsubsection{Operadores especiales}

\begin{description}
  \item[\texttt{Entre \ldots Y \ldots}] Rango de valores. Ej.: \texttt{[Precio] Entre 100 Y 500}
  \item[\texttt{In (\ldots)}] Pertenencia a lista. Ej.: \texttt{[Ciudad] In ("Madrid";"Valencia";"Sevilla")}
  \item[\texttt{Es Nulo}] Valor nulo. Ej.: \texttt{[Email] Es Nulo}
  \item[\texttt{Como}] Comparación con comodines:
    \begin{itemize}
      \item \texttt{*} --- Cualquier cadena de caracteres.
      \item \texttt{?} --- Un único carácter.
      \item \texttt{\#} --- Un único dígito.
      \item \texttt{[lista]} --- Cualquier carácter de la lista.
    \end{itemize}
    Ej.: \texttt{[Nombre] Como "Mar*"} (nombres que empiezan por ``Mar'').
\end{description}

\subsubsection{Operadores lógicos}

\begin{itemize}
  \item Condiciones en la \textbf{misma fila} $\rightarrow$ unidas por \textbf{Y (AND)}.
  \item Condiciones en \textbf{filas distintas} (fila O:) $\rightarrow$ unidas por \textbf{O (OR)}.
\end{itemize}

\subsection{Consultas con parámetros}

Permiten que el usuario introduzca el valor de búsqueda al ejecutar la consulta:

\begin{enumerate}
  \item En la fila Criterios, escribir el nombre del parámetro entre corchetes: \texttt{[Introduce la población]}.
  \item Al ejecutar, Access mostrará un cuadro de diálogo pidiendo el valor.
\end{enumerate}

\begin{importante}
El nombre del parámetro debe ir \textbf{entre corchetes} y no coincidir con ningún nombre de campo existente.
\end{importante}

\subsection{Consultas multitabla}

Una consulta puede obtener datos de varias tablas. Si las tablas están relacionadas, aparecerán automáticamente combinadas (línea de unión).

\subsubsection{Composición interna (INNER JOIN)}
Solo devuelve registros que tienen correspondencia en ambas tablas.

\subsubsection{Composición externa (OUTER JOIN)}
También devuelve registros sin correspondencia en la otra tabla. Se configura haciendo doble clic sobre la línea de unión y eligiendo la opción 2 o 3 en \menu{Tipo de combinación}.

% =============================================
\section{Consultas de resumen}
% =============================================

\subsection{Definición}

Las \textbf{consultas de resumen} agrupan registros y calculan estadísticas. Se activan con el botón \textbf{Totales} en la pestaña Diseño, que añade la fila \textbf{Total:} a la cuadrícula QBE.

\subsection{Funciones de agregado}

\begin{center}
\renewcommand{\arraystretch}{1.3}
\begin{tabular}{>{\bfseries\color{darkblue}}p{2.5cm} p{9cm}}
\toprule
Función & Descripción \\
\midrule
Suma & Suma de valores numéricos. \\
Promedio & Media aritmética. \\
DesvEst & Desviación estándar. \\
Var & Varianza. \\
Mín & Valor mínimo (numérico, texto o fecha). \\
Max & Valor máximo. \\
Primero & Primer registro del grupo (orden cronológico de entrada). \\
Último & Último registro del grupo. \\
Cuenta & Cuenta de valores no nulos. Para contar todos los registros: \texttt{Cuenta(*)}. \\
AgruparPor & Define la columna de agrupación. \\
Expresión & Permite combinar funciones de agregado en una expresión. \\
Dónde & Aplica un criterio de búsqueda \emph{antes} de calcular. \\
\bottomrule
\end{tabular}
\end{center}

\begin{nota}
Los valores \texttt{NULL} no se cuentan en las funciones de agregado. El valor 0 sí se cuenta.
\end{nota}

\subsection{Agrupación de registros}

Usando \textbf{AgruparPor}, Access forma grupos con los registros que tienen el mismo valor en el campo de agrupación. Cada grupo genera una fila en el resultado.

\textit{Ejemplo}: Contar alumnos por población:
\begin{center}
\begin{tabular}{|c|c|}
\hline
\textbf{Campo} & Población | Alumno \\
\hline
\textbf{Total} & AgruparPor | Cuenta \\
\hline
\end{tabular}
\end{center}

% =============================================
\section{Consultas de referencias cruzadas}
% =============================================

Una \textbf{consulta de referencias cruzadas} representa datos de resumen en formato de tabla de doble entrada, con una columna de agrupación como filas y otra como columnas.

Se crea con el \textbf{Asistente para consultas de referencias cruzadas} (\menu{Crear > Asistente para Consultas > Asist. consultas de tabla ref. cruzadas}):

\begin{enumerate}
  \item Elegir la tabla u origen.
  \item Seleccionar el campo(s) de \textbf{encabezado de filas} (hasta 3).
  \item Seleccionar el campo de \textbf{encabezado de columnas} (1 campo).
  \item Elegir el \textbf{campo valor} y la función de agregado.
  \item Opcionalmente, añadir suma de filas.
  \item Asignar nombre y finalizar.
\end{enumerate}

% =============================================
\section{Consultas de acción}
% =============================================

Las \textbf{consultas de acción} modifican los datos. Antes de ejecutarse, Access pide confirmación.

\begin{importante}
Las consultas de acción son \textbf{irreversibles} (salvo deshacer inmediato). Usa siempre la Vista Hoja de datos para revisar qué registros afectará antes de ejecutar.
\end{importante}

\subsection{Consulta de creación de tabla}

Guarda el resultado de una consulta de selección en una \textbf{nueva tabla}.

\begin{enumerate}
  \item Diseñar la consulta de selección.
  \item \menu{Diseño > Tipo de consulta > Crear tabla}.
  \item Indicar el nombre de la nueva tabla.
  \item Ejecutar.
\end{enumerate}

\begin{nota}
Los campos de la nueva tabla heredan el tipo de datos pero \textbf{no} la clave principal ni los índices.
\end{nota}

\subsection{Consulta de actualización}

Modifica los valores de uno o más campos en los registros que cumplen el criterio.

\begin{enumerate}
  \item \menu{Diseño > Tipo de consulta > Actualizar}.
  \item Aparece la fila \textbf{Actualizar a:} en la cuadrícula.
  \item Escribir la expresión del nuevo valor (campo, valor fijo o expresión).
\end{enumerate}

Ejemplo: subir un 10\% el precio: \texttt{[Precio] * 1,1}

\subsection{Consulta de datos anexados}

Añade registros de una tabla de origen al final de una tabla destino.

\begin{enumerate}
  \item \menu{Diseño > Tipo de consulta > Anexar}.
  \item Indicar la tabla destino.
  \item En la fila \textbf{Anexar a:}, indicar el campo destino correspondiente.
\end{enumerate}

\subsection{Consulta de eliminación}

Borra los registros que cumplen el criterio de búsqueda.

\begin{enumerate}
  \item \menu{Diseño > Tipo de consulta > Eliminar}.
  \item Aparece la fila \textbf{Eliminar:}.
  \item Definir el criterio de búsqueda.
\end{enumerate}

\begin{importante}
Si no se especifica ningún criterio, \textbf{se borran TODOS los registros} de la tabla. Los datos borrados \textbf{no se pueden recuperar}.
\end{importante}

% =============================================
\section{Formularios}
% =============================================

\subsection{¿Para qué sirven los formularios?}

Los \textbf{formularios} proporcionan una interfaz gráfica amigable para introducir, visualizar y modificar datos de tablas o consultas.

\subsection{Crear un formulario}

Desde \menu{Crear > Formularios}:

\begin{description}
  \item[Formulario] Crea automáticamente un formulario con todos los campos de la tabla/consulta seleccionada en el Panel de Navegación.
  \item[Diseño del formulario] Formulario en blanco en Vista Diseño para construir manualmente.
  \item[Formulario en blanco] Formulario vacío en Vista Presentación; permite arrastrar campos desde el Panel.
  \item[Asistente para formularios] Guía paso a paso (recomendado).
  \item[Navegación] Formulario tipo menú de página web.
  \item[Más formularios] Varios elementos, Hoja de datos, Formulario dividido, Cuadro de diálogo modal.
\end{description}

\subsection{Asistente para formularios}

\begin{enumerate}
  \item Elegir tabla/consulta origen y campos a incluir.
  \item Seleccionar distribución: \emph{En columnas}, \emph{Tabular}, \emph{Hoja de datos}, \emph{Justificado}.
  \item Asignar título (también será el nombre del formulario).
  \item Elegir: \emph{Abrir para ver datos} o \emph{Modificar el diseño}.
\end{enumerate}

\subsection{Vista Diseño del formulario}

La Vista Diseño tiene tres secciones:

\begin{description}
  \item[Encabezado de formulario] Aparece una sola vez al inicio (títulos, logos).
  \item[Detalle] Muestra los registros del origen (se repite por cada registro).
  \item[Pie de formulario] Aparece una sola vez al final (totales, botones de navegación).
\end{description}

\subsection{Controles}

Un \textbf{control} es cualquier elemento visual del formulario: cuadros de texto, etiquetas, botones, imágenes, listas desplegables, casillas de verificación...

\subsubsection{Tipos de controles principales}

\begin{itemize}
  \item \textbf{Cuadro de texto}: muestra y edita datos de un campo.
  \item \textbf{Etiqueta}: texto fijo descriptivo.
  \item \textbf{Botón de comando}: ejecuta una acción (abrir formulario, ejecutar macro...).
  \item \textbf{Cuadro combinado / Lista}: seleccionar de un conjunto de valores.
  \item \textbf{Casilla de verificación}: campo Sí/No.
  \item \textbf{Imagen}: imagen decorativa o de fondo.
  \item \textbf{Subformulario}: formulario anidado para mostrar datos relacionados (relación 1:N).
\end{itemize}

\subsubsection{Trabajar con controles}

\begin{itemize}
  \item \textbf{Seleccionar}: clic sobre el control. Para varios: \tecla{Ctrl}+clic o arrastrar un rectángulo de selección.
  \item \textbf{Mover}: arrastrar (mueve control + etiqueta). Arrastrar desde el cuadrado gris superior izquierdo mueve solo uno.
  \item \textbf{Cambiar tamaño}: arrastrar los controladores de tamaño (cuadrados pequeños).
  \item \textbf{Copiar/Pegar}: \tecla{Ctrl+C} / \tecla{Ctrl+V}.
  \item \textbf{Eliminar}: \tecla{Supr}.
\end{itemize}

\subsection{Subformularios}

Los \textbf{subformularios} muestran registros relacionados (tabla secundaria) dentro del formulario principal. Son ideales para relaciones \textbf{uno a varios}.

\textit{Ejemplo}: Formulario Cursos con subformulario Alumnos matriculados.

El formulario principal y el subformulario están vinculados automáticamente por los campos de relación.

\subsection{Propiedades del formulario}

Principales propiedades accesibles desde \menu{Diseño > Hoja de Propiedades}:

\begin{description}
  \item[Pestaña Formato] Título, vista predeterminada (Un único formulario / Formularios continuos / Hoja de datos / Formulario dividido), barras de desplazamiento, botones de navegación...
  \item[Pestaña Datos] Origen del registro, filtro, ordenación, Entrada de datos, Permitir agregar/eliminar/editar...
  \item[Pestaña Otras] Emergente, Modal, Ciclo de TAB...
\end{description}

% =============================================
\section{Resumen de atajos de teclado}
% =============================================

\begin{center}
\renewcommand{\arraystretch}{1.3}
\begin{tabular}{>{\ttfamily}p{4cm} p{8cm}}
\toprule
\textnormal{\textbf{Atajo}} & \textbf{Acción} \\
\midrule
Ctrl+G & Guardar \\
Ctrl+Z & Deshacer \\
Ctrl+A & Abrir base de datos \\
Ctrl+B & Buscar \\
Ctrl+C & Copiar \\
Ctrl+V & Pegar \\
Ctrl+X & Cortar \\
Alt & Activar acceso por teclado en la cinta \\
F4 & Desplegar lista en campo \\
F5 & Ir a número de registro \\
Supr & Eliminar selección \\
\bottomrule
\end{tabular}
\end{center}

\vfill
\begin{center}
{\color{accessred}\rule{0.5\linewidth}{1pt}}\\[0.3cm]
{\small\color{gray} UT6 -- Bases de Datos Ofimáticas con Microsoft Access \quad|\quad AOF}
\end{center}

\end{document}
